\documentclass[14pt,a4paper]{extarticle}
\usepackage[T2A]{fontenc}\usepackage[utf8]{inputenc}\usepackage[russian]{babel}
\usepackage{geometry}\geometry{left=30mm,right=15mm,top=20mm,bottom=20mm}
\setlength{\parindent}{1.25cm}\setlength{\parskip}{0.3em}\sloppy
\begin{document}\begin{titlepage}\begin{center}
САНКТ-ПЕТЕРБУРГСКИЙ НАЦИОНАЛЬНЫЙ\\ ИССЛЕДОВАТЕЛЬСКИЙ УНИВЕРСИТЕТ ИТМО
\vspace{1.5cm}

Дисциплина: Бэк-энд разработка
\vspace{1.5cm}

\textbf{Отчёт}\\[0.5cm]
Лабораторная работа 8\\[0.5cm]
\textbf{Нагрузочное тестирование и автомасштабирование}
\vfill
\begin{flushright}Выполнил:\\ Цатинян Артём\\ Группа БР1.2\\[0.8cm] Проверил:\\ Добряков Д. И.\end{flushright}
\vfill Санкт-Петербург\\ 2026 г.\end{center}\end{titlepage}
\section{Задача и состояние}
Провести нагрузочные испытания облачного приложения, исследовать метрики и настроить HPA. В данном документе приведён план: облачный стенд, испытания и настройки HPA пока не подтверждены.
\section{Подготовка}
Использовать проверенный стенд из ЛР7 и сценарии ДЗ10. Подготовить тестовые данные, установить лимиты нагрузки и убедиться, что Prometheus и Grafana получают метрики приложений, узлов и базы. Текущая локальная Grafana из ЛР6 подтверждает мониторинг локального стенда, но не облачного кластера.
\section{Проведение эксперимента}
Зафиксировать число реплик, выделенные ресурсы и начальное состояние. Выполнить прогрев, базовую, повышенную и длительную нагрузку, отслеживая RPS, p95/p99, ошибки и насыщение CPU, памяти, пула подключений и базы. Сохранить графики, параметры и результаты каждой фазы.
\section{Автомасштабирование}
Предлагается настроить Kubernetes Horizontal Pod Autoscaler по CPU, указать минимальное и максимальное число реплик и requests для контейнеров. Требуется работающий источник метрик. Под нагрузкой проверить изменение числа реплик и доступность приложения; отдельно оценить задержку масштабирования и ограничение по узлам.
\section{Оптимизация}
На основании измерений скорректировать requests и limits, размер пула БД и настройки приложения, затем повторить испытание при том же профиле. Сравнение без фактических измерений недостоверно; результаты не подменяются расчётными значениями.
\section{Вывод}
Подготовлена методика проверки производительности и HPA. Фактическое развёртывание, замеры и настройка остаются открытыми.
\end{document}
